\documentclass[11pt]{article}

\usepackage[margin=1in]{geometry}
\usepackage{graphicx}
\usepackage{amsmath}
\usepackage{siunitx}
\usepackage{caption}
\usepackage{fancyhdr}
\usepackage{float}
\usepackage{hyperref}

\pagestyle{fancy}
\setlength{\headheight}{14pt}
\fancyhf{}
\lhead{Project 3: AC Load Measurement}
\rhead{\thepage}
\renewcommand{\headrulewidth}{0.4pt}

\title{\vspace{-2cm}Linear Circuits 2 \\ AC Load Measurement Lab Report}
\author{Group Members: \\ Juvon Mondesir \\ Kenny Gallmon}
\date{}

\begin{document}

\maketitle
\thispagestyle{fancy}

\section*{Objectives}

For this project, the objective was to determine the impedance and power
factor of a hidden AC load at a specific frequency of the V peak to peak, as
the current of the load using a resistor due to the accessibility of the
black box being only wires. Then, we were tasked to find the phase angle and
determine if it was inductive or capacitive.

We were given the following circuit (abstractly):

\begin{figure}[H]
    \centering
    \includegraphics[width=0.5\textwidth]{figures/fig01_zl_circuit.png}
    \caption{Abstract representation of the hidden AC load $Z_L$. Node A is
    connected to a red wire; Node B is connected to a black wire.}
    \label{fig:zl_circuit}
\end{figure}

\section*{Experimental Plan}

To start off, we found the current of the whole series circuit. This
involved putting our multimeter in series with node B and ground while
setting the multimeter to current mode. We pursued finding the voltages of
each element and to do that we utilized a test resistor. After this, we then
proceeded to use Ohm's Law to find the impedance of both the resistor and
the black box. We then need to find the phase difference to find out if the
box is inductive or capacitive in nature, which will tell us what components
are inside the box. To find the phase angle, we used the oscilloscope to
find it with the reference being the original input AC voltage.

\section*{Experimental Results}

Now for the experimental results, we measured the voltage across the shunt
resistor:

\begin{figure}[H]
    \centering
    \includegraphics[width=0.55\textwidth]{figures/fig02_scope_voltage.png}
    \caption{Measured voltage across the shunt resistor, $16.454$\,mV}
\end{figure}

For the current, we measured the current leaving the black box and entering
the same shunt resistor:

\begin{figure}[H]
    \centering
    \includegraphics[width=0.55\textwidth]{figures/fig03_scope_current.png}
    \caption{Measured current leaving the black box, $0.0826$\,mA}
\end{figure}

Using these measurements, some hand calculations were made:

\begin{figure}[H]
    \centering
    \includegraphics[width=0.7\textwidth]{figures/fig04_handcalc.png}
    \caption{Hand calculations for impedance, phase difference, and power
    factor}
\end{figure}

\[
Z = 6.1\,\text{k}\Omega + j0
\]

\section*{Results Comparison}

\begin{figure}[H]
    \centering
    \includegraphics[width=0.95\textwidth]{figures/fig05_scope_waveform.png}
    \caption{Voltage (C2) and current (C3) waveform}
\end{figure}

For this section, after finding out that our phase angle is zero, we
determined that our black box is not capacitive or inductive in nature.
Thus, this means our mystery circuit is purely resistive, aka an RR circuit.
The power factor is 1, which means physically, the voltage and current are
perfectly in phase, which can be seen in the image above.

\section*{Conclusion}

In this experiment, the objective was to identify the characteristics of the
unknown AC load by determining its impedance, power factor, and its nature.
This is done by measuring the voltage and current across the concealed
circuit and a test resistor, allowing us to analyze the circuit's behavior.
Our measurements revealed an impedance of $6.1$\,k$\Omega$ and a phase angle
of zero, which led to the conclusion of the concealed circuit being purely
resistive.

\end{document}